\section{Serialization changes the learning problem}
\label{sec:serialization}

\subsection{A bijection and its accounting}

For an even $s\times s$ lattice, write $P=(s/2)^2$ for the number of
$2\times2$ patches. For alphabet size $q$, the encoding
\begin{equation}
 z_{t,p}=\beta+\sum_{a=0}^{1}\sum_{b=0}^{1}
        q^{2a+b}X^t_{2r+a,2c+b},
 \quad p=(s/2)r+c
 \label{eq:patch}
\end{equation}
is a bijection between four cell states and $q^4$ patch symbols.
Repeated division by $q$ recovers its four digits, proving invertibility.
$\beta$ is a vocabulary offset; the binary pilot uses $\beta=256$.
Two additional marker tokens make
\begin{equation}
 \ell=P+2,\qquad M_{\rm seq}=T(P+2).
\end{equation}
For $q=2,s=8,T=16$, the values are
\[
 q^4=16,\quad P=16,\quad\ell=18,\quad M_{\rm seq}=288.
\]
For a $12\times12$, ten-state construction, they become
$q^4=10\,000$, $P=36$, and $\ell=38$.
The number of possible patch patterns is not necessarily the allocated
transformer vocabulary. An invertible recoding of a complete trajectory
preserves its entropy; markers and integer storage width change file size,
not its underlying discrete information.

\subsection{Direct-parent visibility}

Index tokens from zero. Patch $p$ in frame $t$ occupies
$j(t,p)=t\ell+1+p$. A source cell in patch $r$ of the preceding frame
is encoded at $j(t-1,r)$, so its backward token lag is
\begin{equation}
 D(p,r)=j(t,p)-j(t-1,r)=\ell+p-r.
 \label{eq:lag}
\end{equation}
Since $0\leq p,r\leq P-1$,
\begin{equation}
 3\leq D(p,r)\leq2P+1.
 \label{eq:lagrange}
\end{equation}
This includes arbitrary rewired one-step parents; there is no assumption
that their patches are spatially adjacent.

\begin{proposition}[Full-window one-step visibility]
If predicting a patch token with a complete preceding-token window of
length $W\geq2P+1$, all of its preceding-frame source patches are
available to the causal predictor, regardless of the rewiring.
\end{proposition}
\begin{proof}
Every required source token precedes the target, and its lag is bounded
by \eqref{eq:lagrange}. The specified window therefore contains it.
\end{proof}

For the binary pilot, $2P+1=33<128$. Thus arbitrary direct parents
\emph{can fit} inside the configured context. This does not show that
the model can identify those parents, extract the required bits, infer
the local rule, or learn it efficiently. For $P=36$, the corresponding
threshold is 73, also below a 1024-token context.

\textbf{Crop-boundary qualification.}
The historical training loss is applied at every position of a randomly
cropped sequence. A target at offset $j\in\{1,\ldots,W\}$ from the first
input has only $j$ preceding input tokens, not $W$ of them. The full-window
proposition guarantees visibility only if $j\geq33$ in the binary case.
For $W=128$, offsets 33 through 128 give $96$ qualifying positions;
the first 32 can lose some direct parents. Which positions are patch
targets depends on the crop's frame offset. We therefore cannot dismiss
all truncation effects merely because the configured context is 128.
A full-frame-aware objective or a sufficiently masked prefix is needed
to make the guarantee apply to every evaluated patch target.

\subsection{An exact invariant of degree-preserving rewiring}

Let $p(i)$ be the patch index of vertex $i$. Choose a vertex and a
non-self relation uniformly and let $J$ be its source. Every patch contains
four vertices. Both endpoints are uniform because each relation is a
permutation. Therefore
\begin{equation}
 \E[p(i)]=\E[p(J)]=\frac{P-1}{2},
 \qquad \boxed{\E[D]=\ell.}
 \label{eq:meanlag}
\end{equation}
This holds for every fixed typed-permutation topology, including grid
and rewired graphs. Rewiring cannot change this mean lag when the
average weights all relation-slot occurrences equally.
It can change the lag distribution and which specific nodes are related.
If the neural rule weights some slots more strongly, an influence-weighted
quantity requires a different definition.

Writing $\sigma_p^2=(P^2-1)/12$,
\begin{equation}
 \Var(D)=2\sigma_p^2-2\Cov(p(i),p(J)).
 \label{eq:lagvariance}
\end{equation}
Thus equal mean lag does not imply equally dispersed dependencies.
This identifies a possible serialization effect without assuming a
language-transfer effect.

\subsection{Exact grid variance for the binary pilot}

On the $8\times8$ grid,
$p(r,c)=4\lfloor r/2\rfloor+\lfloor c/2\rfloor$.
Let $A=\lfloor r/2\rfloor$, $B=\lfloor c/2\rfloor$, with $r,c$ uniform.
For a unit periodic row shift, $\Delta A$ has squared values
$0,1,0,1,0,1,0,9$, hence
\begin{equation}
 \E[\Delta A]=0,\qquad \E[(\Delta A)^2]=\frac{12}{8}=\frac32.
\end{equation}
A column shift has the same properties for $\Delta B$. Among the eight
Moore directions, six shift the row and six shift the column. For a
diagonal shift the row and column increments are independent with mean
zero. Averaging $(4\Delta A+\Delta B)^2$ therefore gives
\begin{align}
 \Var(D_{\rm grid})
 &=16\cdot\frac32\cdot\frac68
       +\frac32\cdot\frac68\\
 &=\frac{153}{8}=19.125,\qquad
 {\rm sd}(D_{\rm grid})=4.373213921.
 \label{eq:gridlagvariance}
\end{align}
The result includes periodic wrap-around; dropping the wrap terms would
give a different and incorrect calculation.

For comparison, consider an idealized vertex-exchangeable non-self
source, uniform among the other 63 vertices. For independent uniform
vertices with replacement, $2\sigma_p^2=2(255/12)=42.5$.
The excluded self pairs contribute zero to the squared difference.
Conditioning on distinct vertices thus gives
\begin{equation}
 \begin{aligned}
 \Var(D_{\rm uniform\ nonself})
 &=\frac{64}{63}\cdot42.5=\frac{2720}{63}
 =43.17460317,\\
 {\rm sd}(D_{\rm uniform\ nonself})&=6.570738404.
 \end{aligned}
 \label{eq:randomlagvariance}
\end{equation}
This comparison law is not a calculation of the pilot's partial-swap
graph, whose endpoints are not assumed to follow this idealized law.
It shows how lag dispersion can rise while mean lag remains exactly 18.
Neither \eqref{eq:gridlagvariance} nor \eqref{eq:randomlagvariance}
establishes that dispersion is beneficial.

